\chapter{הגדרת הבעיה}\label{ch:problem}

\section{משפט הבעיה}
\begin{keybox}[משפט הבעיה]
לעורכי חדשות אין דרך מהירה, אמינה ומתועדת לקבוע לפני פרסום אם תמונה היא תיעוד אמיתי או תוצר של בינה מלאכותית. לכן תמונות שקריות מתפרסמות תחת מותג אמין, הציבור מוטעה, ומערכת החדשות מאבדת אמון שקשה מאוד לשקם.
\end{keybox}

לבעיה שני רבדים שצריך לפתור יחד:
\begin{enumerate}
\item \textbf{הרובד הטכני, אי ודאות:} עורך שמקבל תמונה ממקור חיצוני (רשת חברתית, קורא, סוכנות קטנה) צריך להחליט תוך דקות. כלי הזיהוי הזמינים נותנים תשובות סותרות, אינם מסבירים את עצמם, ומאבדים דיוק בתנאי שטח ואחרי עיבוד של התמונה~\cite{ren2025}. רוב התמונות אינן נושאות תעודת מקור~\cite{c2pa2026}.
\item \textbf{הרובד האתי, תמריץ כלכלי:} גם כשיש ספק, הלחץ לפרסם ראשון ולמשוך טראפיק דוחף לפרסום. אין כלל מוסכם שקובע מה עושים במצב ״לא ודאי״, איך מתייגים, ומי אחראי. כללי אתיקה קיימים מנוסחים ברמת עיקרון (״אל תעוות מידע חזותי״~\cite{spj2014}) ולא ברמת נוהל עבודה.
\end{enumerate}

\section{בעלי עניין}
איור~\ref{fig:stake} ממפה את בעלי העניין לפי מידת ההשפעה שלהם על הפתרון ומידת הפגיעה שהבעיה גורמת להם. המיפוי קובע את סדר הפנייה: עורכים ומנהלי מערכות הם הלקוח המשלם, רשתות הפרסום הן מנוף ההשפעה, והציבור הוא הנפגע שבשמו הפתרון פועל.

\begin{figure}[!htb]
\begin{ltrdiagram}
% stake: TikZ / pgfplots source, built inside ltrdiagram (LTR)
\begin{tikzpicture}[font=\sffamily\small,
  pt/.style={circle,fill=acc2,inner sep=2.2pt},
  lb/.style={font=\sffamily\scriptsize,align=center}]
\draw[->,thick,ink] (0,0)--(10.4,0) node[below left,lb]{\he{השפעה על הפתרון} $\rightarrow$};
\draw[->,thick,ink] (0,0)--(0,6.4) node[above,lb]{\he{פגיעה מהבעיה} $\uparrow$};
\draw[dashed,gray] (5.2,0)--(5.2,6.2);
\draw[dashed,gray] (0,3.1)--(10.2,3.1);
\node[lb,gray] at (2.6,6.0) {\he{לשתף ולהגן}};
\node[lb,gray] at (7.8,6.0) {\he{שותפים מרכזיים}};
\node[lb,gray] at (2.6,0.35) {\he{לעדכן}};
\node[lb,gray] at (7.8,0.35) {\he{לרתום}};
\foreach \x/\y/\t in {
 1.4/5.2/{\he{הציבור / הקוראים}},
 3.6/4.4/{\he{צלמי עיתונות}},
 8.6/5.0/{\he{עורכים ומנהלי מערכות}},
 7.6/4.3/{\he{מועצת העיתונות / גופי אתיקה}},
 9.2/1.9/{\he{רשתות פרסום פרוגרמטי}},
 7.0/1.1/{\he{פלטפורמות רשת חברתית}},
 3.0/1.6/{\he{רגולטורים}},
 6.9/2.5/{\he{יצרני מצלמות ו-\en{C2PA}}}}{
 \node[pt] at (\x,\y) {};
 \node[lb,anchor=south] at (\x,\y+0.08) {\t};}
\end{tikzpicture}

\end{ltrdiagram}
\caption{בעלי עניין: העורכים משלמים, המפרסמים הם המנוף, והציבור הוא הנפגע}\label{fig:stake}
\end{figure}

האיור מראה מי משפיע על מי. טבלה~\ref{tab:stake} מפרטת לכל בעל עניין מה הוא צריך ומה הפרויקט נותן לו.


\begin{table}[!htb]\centering
\caption{בעלי העניין, הצורך שלהם ומה הפרויקט נותן להם}\label{tab:stake}
\begin{tabularx}{\linewidth}{@{}l Y Y@{}}
\toprule
\textbf{בעל עניין} & \textbf{הצורך / הכאב} & \textbf{מה \en{VeriLens} נותנת} \\
\midrule
עורך משמרת & החלטה תוך דקות בלי לטעות & ציון מכויל, הסבר חזותי ותיעוד ההחלטה \\
מנהל/ת מערכת & הגנה על המוניטין ועמידה ברגולציה & לוח בקרה, דוח שקיפות והסמכת \en{TRUTH-IMG} \\
צלם/ת עיתונות & הוכחה שהצילום שלו/ה אמיתי & חתימת \en{C2PA} באפליקציה ושמירת התעודה עד הפרסום \\
קורא/ת & לדעת מה אמיתי & תווית גלויה ותעודת מקור בלחיצה \\
מפרסם / רשת פרסום & לא לממן שקר (\en{Brand Safety}) & ממשק \en{API} ״ציון שלמות״ לאתר \\
רגולטור & אכיפת סעיף \num{50} & תיעוד ביקורת (\en{Audit Trail}) לפי דרישה \\
\bottomrule
\end{tabularx}
\end{table}

\section{גבולות הפרויקט}

כדי שהחוזה יהיה ברור, טבלה~\ref{tab:scope} מגדירה מה נכלל בפרויקט ומה נשאר מחוצה לו.

\begin{table}[!htb]\centering
\caption{גבולות התחום: בתוך הפרויקט ומחוצה לו}\label{tab:scope}
\begin{tabularx}{\linewidth}{@{}Y Y@{}}
\toprule
\textbf{בתוך הפרויקט} & \textbf{מחוץ לפרויקט (במפורש)} \\
\midrule
תמונות סטילס (\en{JPEG}, \en{PNG}, \en{WebP}, \en{HEIC}) & וידאו ואודיו (מתוכנן לשלב \num{2}, פרק~\ref{ch:business}) \\
מערכות חדשות דיגיטליות ומחלקות דסק & אכיפה מול משתמשים פרטיים ברשתות \\
אימות לפני פרסום, וסריקה של ארכיון \num{90} הימים האחרונים & הסרת תוכן מאתרים של צד שלישי \\
עברית, אנגלית, ערבית וצרפתית בממשק העורך & ממשק לשפות נוספות \\
התקן \en{TRUTH-IMG} ותוכנית הסמכה & חקיקה או ייעוץ משפטי אישי \\
\bottomrule
\end{tabularx}
\end{table}

\section{הנחות יסוד ותלויות}
ההצעה נשענת על ההנחות שבטבלה~\ref{tab:assume}. אם אחת מהן לא תתקיים, הלו״ז, התקציב או היעדים עלולים להשתנות, ולכן לכל הנחה מוגדר מה קורה אם היא נופלת. זהו נוהג מקובל בהצעות מחיר ובאומדני עלות: מנהל שאינו מכיר את ההנחות אינו יכול לשפוט את האומדן~\cite{gao2020}.

\begin{table}[!htb]\centering
\caption{הנחות יסוד, תלויות, ומה קורה אם הן לא מתקיימות}\label{tab:assume}
\small
\begin{tabularx}{\linewidth}{@{}l Y Y@{}}
\toprule
\textbf{מזהה} & \textbf{הנחה או תלות} & \textbf{אם היא לא מתקיימת} \\
\midrule
\req{A-01} & חמש מערכות חדשות יחתמו על הסכם פיילוט עד חודש \num{10} & הפיילוט נדחה; לכן פונים לשבע מועמדות עבור חמישה מקומות \\
\req{A-02} & סוכנויות צילום ירשו \num{600} אלף צילומים לאימון במחיר של עד \num{0.25} ₪ לצילום & עלייה בתקציב, שנבלעת ברזרבה עד \num{0.4} ₪ לצילום \\
\req{A-03} & ספקי חיפוש תמונה הפוך יאפשרו גישה בממשק בתשלום & שכבת ההקשר תיבנה על מאגר עצמי, בעלות של כ-\num{3} חודשי עבודה \\
\req{A-04} & תקן \en{C2PA} יישאר בגרסה תואמת לאורך הפרויקט & עדכון הרכיב בשכבה \num{1}, כחודש עבודה \\
\req{A-05} & הנחיות סעיף \num{50} האירופי לא ישתנו מהותית עד חודש \num{18} & עדכון תוויות ונהלים; מעקב משפטי רבעוני \\
\req{A-06} & המרצה יאשר את מאגר \en{VL-Wild} כאמת המידה לבדיקה & נדרשת בקשת שינוי לפני חודש \num{5} \\
\bottomrule
\end{tabularx}
\end{table}

\subsection*{מה ייחשב פתרון}
הבעיה תיחשב פתורה במערכות הפיילוט אם יתקיימו שלושה תנאים: (\num{1}) שיעור התמונות הסינתטיות שמתפרסמות כתיעוד יורד ב-\num{80\%} לפחות לעומת קו בסיס שנמדד בחודש הראשון; (\num{2}) זמן ההחלטה החציוני של עורך אינו עולה על \num{15} דקות; (\num{3}) לכל תמונה שפורסמה קיימת רשומת החלטה מלאה ביומן. היעדים המספריים המלאים מפורטים בפרק~\ref{ch:goals}.
